\subsection*{トピックと参考文献の対応}
\addcontentsline{toc}{subsection}{8 トピックと参考文献の対応}
この表は、SWEBOK が示す「トピック対参考資料の行列」を、学習用に簡略化したものである。詳しく学ぶときは、各トピックに対応する文献を読む。

\par\smallskip\noindent\refstepcounter{table}\label{tab:ch08-08}表 \thetable: トピックと参考文献の対応におけるトピック・主な参考資料\par\smallskip
表 \thetable は、トピックと参考文献の対応におけるトピック・主な参考資料を比較・確認しやすい形で整理したものである。\par\smallskip
\begin{longtable}{p{0.30\linewidth}p{0.64\linewidth}}\toprule
トピック & 主な参考資料 \\ \midrule
\endfirsthead\toprule
トピック & 主な参考資料 \\ \midrule
\endhead
1. 構成管理プロセスの管理 & IEEE Std 828; Hass; Sommerville \\
1.1 組織における位置づけ & IEEE Std 828; Hass; Sommerville \\
1.2 制約と指針 & IEEE Std 828; Hass; Moore \\
1.3 構成管理の計画 & IEEE Std 828; Hass; Sommerville \\
1.3.3 ツール選定と実装 & Hass \\
1.4 ソフトウェア構成管理計画 & IEEE Std 828; Hass \\
1.5 構成管理の監視 & Hass \\
2. ソフトウェア構成識別 & IEEE Std 828; Hass \\
2.1 管理対象品目の識別 & IEEE Std 828; Hass \\
2.1.1 ソフトウェア構成 & ISO/IEC/IEEE 24765; IEEE Std 828 \\
2.1.2 ソフトウェア構成品目 & ISO/IEC/IEEE 24765; Hass \\
2.2 識別子と属性 & IEEE Std 828; Hass \\
2.5 関係方式の定義 & Hass \\
2.6 ソフトウェアライブラリ & IEEE Std 828; Hass \\
3. ソフトウェア構成変更統制 & IEEE Std 828; Hass; Sommerville; Moore \\
3.1 変更の要求・評価・承認 & Hass; Sommerville; Moore \\
3.2 ソフトウェア変更の実施 & Sommerville; Moore \\
4. 構成状態記録・報告 & IEEE Std 828; Hass; Moore \\
5. ソフトウェア構成監査 & IEEE Std 828; Moore \\
6. リリース管理と提供 & IEEE Std 828; Hass; Sommerville \\
7. 構成管理ツール & Hass \\
\bottomrule\end{longtable}
